\section{Discussion}
\label{sec:discussion}

The results illustrate why high-frequency digital trace data can add information beyond aggregate estimates of shutdown costs. A national GDP series cannot show whether digital workers cease production, shift to offline tasks, or accumulate work that is released later. GitHub events and metadata measure observed public activity and delayed publication; they cannot observe work that never becomes public or distinguish all unobserved adaptation from lost production.

The main estimate---an 86.8 percent decline in the daily active-user rate relative to the synthetic counterfactual (\cref{tab:sc_main})---is best interpreted as a collapse in observable online software-production activity. The mechanism shows a small observed backlog alongside a collapse in author-time levels: creation of code later observed publicly falls roughly 96--97 percent from the pre-shutdown comparison, and backlog represents about 2.5 percent of normal volume. The original 28.27 percent backlog share conditions on a much smaller calendar-cohort denominator. It does not establish that missing activity was displaced across time, or quantify destruction of work that never became public.

At the same time, delayed synchronization cannot account for the entire collapse. Most eligible users record no productive public activity during the shutdown, and the recovery backlog is much smaller than the counterfactual volume implied by the active-rate and shipping estimates. Connectivity therefore appears to matter not only for publishing already-written code but also for complementary tasks: accessing remote dependencies, coordinating with collaborators, reviewing code, resolving issues, merging pull requests, and releasing software. The especially large declines in shipping and collaboration are consistent with these channels.

The sectoral results add a dynamic dimension. In aggregate, Bangladesh's gap relative to the donor mean returned within its pre-shutdown range in the first week after restoration (\cref{fig:did_es}), but that wide and volatile range is a weak recovery criterion. The separate daily, rolling-mean and cumulative measures in Section~\ref{sec:results} give different return dates and do not show full compensation by September 4. Reconnection did not restore observable activity at once: weeks after the shutdown, activity in Education-, Health-, and especially Accommodation/Food-related repositories remained below its pre-shutdown level, while Finance- and Information-related activity recovered within about three weeks. This pattern is consistent with the costs of a shutdown extending beyond the days without connectivity and falling unevenly across types of projects. The weaker recovery of pre-existing repositories, relative to sector-wide activity, further suggests that part of the recovery took the form of new projects or switching across repositories, which an analysis restricted to pre-existing projects would miss. Because the sectoral analysis has no untreated comparison group and classifies repositories rather than workers, these interpretations remain suggestive.

The national cumulative comparison is more informative about recovery than a first return inside a wide pre-period range. Relative to the unchanged synthetic series, 8,145 active user-days are missing in the six-day shutdown window; later positive gaps partly offset, but never fully repay, the cumulative shortfall, which is 3,633 user-days at September 4. Neither this arithmetic offset nor sectoral baseline crossings identify a causal rebound. The data do not evaluate particular resilience policies, so no comparative policy recommendation is drawn.
